\clearpage
\section{Описание применения генеративных моделей}
\label{sec:ai-usage}

При подготовке работы использовался программный агент OpenAI Codex компании OpenAI. На раннем этапе переработки применялась модель GPT-5; при финальной критической вычитке и исправлении исследования --- GPT-6. Работа с агентом велась через клиент Codex с доступом к файлам проекта и возможностью запуска программ анализа. Описание продукта доступно по адресу \url{https://openai.com/codex/}. Период использования при переработке работы --- с 12 августа по 18 сентября 2026 года.

Генеративные модели участвовали в содержательной и технической переработке исследования. Их применение охватывало следующие задачи:
\begin{enumerate}
    \item критический разбор первоначального дизайна исследования и выявление расхождений между текстом, программным кодом и данными;
    \item подготовку вариантов структуры разделов, уточнение формулировок и научно-деловое редактирование текста;
    \item рефакторинг программ подготовки данных, разведочного анализа и прогнозного эксперимента;
    \item разработку и редактирование программ построения графиков;
    \item поиск ошибок в обработке дат набора, сезонов, пропусков и последующих статусов;
    \item подготовку автоматических тестов, запуск вычислений и проверку сборки документа в системе \LaTeX.
    \item на финальном этапе --- критическое чтение текста курсовой работы, анализ последовательности исследования и выявление расхождений между постановкой, программным кодом, расчётами и выводами.
\end{enumerate}

Модель получала запросы на русском языке с указанием изменяемых файлов, требований к исследованию и замечаний научного руководителя. Сгенерированные предложения не принимались автоматически. Численные утверждения сверялись с сохранёнными Таблицами и отчётами, изменения программного кода проверялись тестами, а ссылки на литературу --- по библиографическим данным и первичным публикациям. После внесения изменений повторно запускались подготовка данных, построение графиков, тесты и сборка документа.

Использование модели ускорило поиск несогласованностей и техническую переработку проекта, однако потребовало обязательной проверки. В ходе работы модель предлагала чрезмерно уверенные интерпретации описательных связей, допускала неудачные формулировки и не всегда правильно учитывала контекст всей работы. Такие предложения исправлялись или исключались. Генеративная модель не использовалась как источник эмпирических данных и не заменяла расчёты программ анализа.

Применение агента помогло обнаружить почти совпадающие групповые повторы, различия в агрегировании метрик и несогласованные интерпретации. 

Автор определял постановку исследования, принимал решения о включении методов и результатов, проверял предложенные изменения и несёт полную ответственность за содержание работы. Описание использования модели составлено в соответствии с требованием НИУ ВШЭ раскрывать цели, конкретную технологию, способ и результат её применения~\cite{hse-ai-guidelines}.
